\subsection{Integral information in the formal-diffeomorphism question}
\label{dev:1977-jets}

\paragraph{The question and its scope.}
The last line of \GSpage{P000773}, independently checked against the image, asks
\ru{посчитать когомологии группы} \(\operatorname{Diff}_1(\mathbb Z)\).
The preceding page defines the inverse limit of jet groups and the displayed
series \(t+a_1t^2+a_2t^3+\cdots\). Neither the coefficient module nor the
cochain theory is specified. The strongest literal formulation is therefore a
request for a computation, with these two parameters still to be supplied.
It is not an assertion that discrete group cohomology, continuous cohomology,
and Lie-algebra cohomology coincide.

For a bounded development we explicitly choose the integral counterpart of
the finite-jet interpretation. Put
\[
 J=\{t+a_1t^2+a_2t^3+\cdots:a_i\in\mathbb Z\},\qquad
 J_N=J/\{f:f(t)\equiv t\pmod {t^{N+1}}\}\quad(N\geq2).
\]
Multiplication is composition, \(J\) has its coefficientwise topology, and
\(J_N\) is discrete. For a trivial discrete coefficient group \(A\), let
\[
 C^r_{\mathrm{jet}}(J;A)
   =\bigcup_N\operatorname{Map}(J_N^r,A),\qquad
 H^r_{\mathrm{jet}}(J;A)=H^r(C^\bullet_{\mathrm{jet}}(J;A)).
\]
Inflation makes these unions literal subgroups of the cochains on \(J\).
Every cochain in this definition depends on a common finite jet of all its
arguments. This is an editorial choice supported by the preceding inverse-limit
discussion, not a coefficient convention recoverable from the question itself.

\paragraph{What the later literature establishes.}
\citet[Theorems~1.2 and~1.5; Proposition~4.4]{babenkobogatyi2008substitutions}
identify the closed-commutator quotient with
\(\mathbb Z^2\oplus(\mathbb Z/2)^2\), and identify finite-jet real cohomology
with the continuous cohomology of \(L_1\). Their manuscript was received on
24 November 2006 and published in 2008. Thus the real version has a direct
later solution. \citet[Introduction and Theorem~1.0.2 of the preprint]{pavezcornejo2025abelianization}
extends the topological abelianization calculation to deeper substitution
subgroups. That work first circulated on 19 July 2024 and appeared in 2025;
its commutator notation explicitly denotes a closure. Neither statement
supplies a computation of the unrestricted discrete integral cohomology ring.
As of 10 November 2026, the sources checked here do not establish that stronger
computation. This is a limit of the verified match, not a proof of its absence
from all literature.

The development below proves two elementary consequences for the chosen
integral cochain theory. They are supplied here for the companion, with no
claim of research priority. They also expose precisely why replacing the
group by its Lie algebra would lose information.

\begin{lemma}[Finite-jet cochains and coefficient sequences]\ours{}
Filtered inflation gives
\(H^r_{\mathrm{jet}}(J;A)=\varinjlim_N H^r(J_N;A)\).
A short exact sequence of discrete trivial coefficient groups induces a
short exact sequence of these cochain complexes and hence its usual long
exact cohomology sequence.
\end{lemma}
\begin{proof}
Any finite list of cochains is inflated from one sufficiently large \(J_N\).
Consequently a cocycle condition or a coboundary equation in the union holds
at a finite stage. This proves the asserted equality of cohomology and the
filtered colimit. If \(B\to C\) is surjective and a \(C\)-valued cochain
factors through \(J_N^r\), choose a lift of each of its values in \(B\).
The resulting lift still factors through \(J_N^r\). Exactness at the other
terms is pointwise. The connecting homomorphisms follow from the ordinary
cochain argument.
\end{proof}

Write the first coefficients as
\[
 f(t)=t+at^2+bt^3+ct^4+dt^5+et^6+ht^7+\cdots.
\]
For clarity \(h\) here is a coefficient, not another group element. The
explicit quotient in \citet[Proposition~2.1]{babenkobogatyi2008substitutions}
has coordinates
\[
 x=a,\qquad y=b-a^2,\qquad
 \varepsilon_1=c+d+\frac{b(b+1)}2\pmod2,\qquad
 \varepsilon_2=h+db+d\pmod2.                 \tag{J1}
\]
The formula involves integer-valued expressions and depends only on the
seventh jet. Its kernel is the closed commutator subgroup by
\citet[Proposition~2.9, pp.~249--250]{babenkobogatyi2008substitutions}.
The literature input in the next proposition is
exactly this theorem; no claim about the algebraic, unclosed commutator is
needed.

\begin{proposition}[First cohomology and all torsion in degree two]\ours{}
With the definitions and trivial coefficients above,
\[
 H^1_{\mathrm{jet}}(J;\mathbb Z)=\mathbb Zx\oplus\mathbb Zy,
 \qquad
 \operatorname{Tor}H^2_{\mathrm{jet}}(J;\mathbb Z)
       \cong(\mathbb Z/2)^2.                 \tag{J2}
\]
The two torsion generators are represented by
\[
 B_i(f,g)=\frac{E_i(f)+E_i(g)-E_i(f\circ g)}2,
 \qquad i=1,2,                             \tag{J3}
\]
where \(E_i\in\{0,1\}\) is the standard lift of \(\varepsilon_i\).
\end{proposition}
\begin{proof}
Degree-one cocycles with trivial coefficients are homomorphisms; the
finite-jet condition is equivalent to continuity to a discrete target.
Their kernels are closed and contain every commutator, hence contain its
closure. Conversely every homomorphism from
\(A_0=\mathbb Z^2\oplus(\mathbb Z/2)^2\) is continuous after composition
with (J1), since that quotient map already factors through a finite jet.
Therefore
\[
 H^1_{\mathrm{jet}}(J;D)=\operatorname{Hom}(A_0,D)
\]
for each discrete trivial group \(D\). Taking \(D=\mathbb Z\) gives the
first assertion.

Apply the lemma to \(0\to\mathbb Z\to\mathbb Q\to\mathbb Q/\mathbb Z\to0\).
The relevant part of the exact sequence is
\[
 \operatorname{Hom}(A_0,\mathbb Q)\longrightarrow
 \operatorname{Hom}(A_0,\mathbb Q/\mathbb Z)\longrightarrow
 H^2_{\mathrm{jet}}(J;\mathbb Z)\longrightarrow
 H^2_{\mathrm{jet}}(J;\mathbb Q).
\]
The first map is surjective on the two free coordinates, while its image
on the two order-two coordinates is zero. Its cokernel is exactly
\((\mathbb Z/2)^2\), which therefore injects into integral degree two as
the kernel of the coefficient map to \(\mathbb Q\). Rational cochains
form a complex of rational vector spaces, so their cohomology is
torsion-free. All integral torsion lies in that kernel; conversely the
kernel just found is finite, so it consists entirely of torsion. This
proves the second assertion, including the absence of odd torsion and of
elements of order four or higher in this degree and cochain theory.

Under \(\operatorname{Hom}(A_0,\mathbb Q/\mathbb Z)\), the torsion
coordinates are the characters \(\varepsilon_i/2\). Their connecting
cocycles are exactly (J3). Additivity modulo two makes (J3) integral, and
\(\delta^2=0\) gives \(\delta B_i=0\). Also \(2B_i=\delta E_i\).
Neither character lifts to a rational homomorphism on the torsion summand,
and they are independent there. Their two classes thus have exact order
two and generate the kernel calculated above.
\end{proof}

\begin{proposition}[An integral primitive for the degree-one product]\ours{}
For \(w(f)=c-3ab+2a^3\), the inhomogeneous coboundary is
\[
 (\delta w)(f,g)=x(f)y(g).                  \tag{J4}
\]
In particular \([x]\smile[y]=0\) over \(\mathbb Z\), both for finite-jet
cochains and in the ordinary cohomology of the underlying discrete group.
The latter assertion concerns these particular classes only.
\end{proposition}
\begin{proof}
For \(g(t)=t+At^2+Bt^3+Ct^4+\cdots\), direct substitution gives
\[
 a(f\circ g)=a+A,\quad
 b(f\circ g)=b+B+2aA,\quad
 c(f\circ g)=c+C+aA^2+2aB+3bA.
\]
Substitute these three identities in
\(\delta w(f,g)=w(g)-w(f\circ g)+w(f)\).
The cubic and mixed terms cancel, leaving \(a(B-A^2)\), as required.
All functions used depend on the fourth jet, so the equation holds in
both cochain complexes claimed.
\end{proof}

There is a second algebraic check. Put
\(z=c-a^3-2a(b-a^2)\). The fourth-jet composition law in the integer
coordinates \((x,y,z)\) is
\[
 (x,y,z)(X,Y,Z)=(x+X,y+Y,z+Z+yX).
\]
The coordinate change has an integer inverse, so this quotient is the
integer Heisenberg group. Since \(w=z-xy\), its displayed law reproduces
(J4), independently of the direct expansion. The degree-one squares also
vanish: for an integer homomorphism \(u\),
\(\delta\binom{u}{2}=-u\smile u\). Together with graded commutativity,
this checks that every product of two classes in the first group of (J2)
is zero.

\paragraph{Checks and their exact force.}
\texttt{code/check\_jets.py} uses symbolic substitution for (J4), a separate
integer-convolution implementation for composition, 1,500 deterministically
seeded triples with coefficients between \(-50\) and \(50\), and exhaustive
residue checks of the two parity characters (1,024 and 4,096 pairs).
All identities, including the two cocycle equations, passed. The proof of
the torsion assertion was checked a second way with
\(0\to\mathbb Z\xrightarrow{\,2\,}\mathbb Z\to\mathbb F_2\to0\):
\(H^1_{\mathrm{jet}}(J;\mathbb F_2)\) has dimension four, the reduction of
integral degree one has dimension two, and the resulting Bockstein image
has dimension two. The rational-coefficient argument is still needed to
exclude higher and odd torsion; the modulo-two check alone cannot do that.
The computational output is included. A separate review re-derived the
cup-product sign and checked the exact-sequence proof against the published
kernel statement; it found no mathematical defect. It did sharpen the
locator from the displayed quotient map to Proposition~2.9, which proves
that its kernel is the closed commutator. These checks concern formulas and
logical implications, not an independent proof of that published classification.

\paragraph{What remains, and a bounded next step.}
These results determine degree one, the torsion in degree two, and its two
explicit cocycles. They do not determine the torsion-free integral part of
degree two, higher integral groups or their products. They also do not
identify finite-jet cohomology with arbitrary continuous or discrete
cohomology. The remaining obstacle is control of integral inflation maps
along the central extensions \(J_{N+1}\to J_N\); rational Lie-algebra
comparison discards precisely that arithmetic information. A further
bounded run could construct integral resolutions for several successive
jet quotients, compute their degree-two inflation maps, and decide which
integral lattices and denominators survive. A finite range would be
evidence for a stabilization assertion, not a proof of one. Nothing in
the notebook or the sources checked establishes why the original speaker
did not pursue the full integral question; the documented topology and
coefficient distinctions explain a mathematical obstacle without inventing
a historical motive.

\subsection{Verification annotation: the Laplace spectrum on the positive Witt algebra}
\label{dev:1977-laplacian-check}
This is a check of a published class-A thread, not a newly selected open
problem. On \GSpage{P000784--P000785}, take
\([e_i,e_j]=(j-i)e_{i+j}\) for \(i,j\geq1\), declare the increasing wedge
monomials orthonormal, and write
\(\Delta_q=d_q^{\mathsf T}d_q+d_{q+1}d_{q+1}^{\mathsf T}\).
Each weight space is finite-dimensional. The corrected spectral statement
is given by \citet[Theorem~2, p.~324]{fuchswilmarth2010laplacian}; the following checks locate a
specific error in the notebook's assertion about distinct eigenvalues.

\begin{example}[A collision of the displayed spectral labels]\oursComputed{}
The two gap-three partitions of weight 25
\[
 (1,10,14),\qquad(3,7,15)
\]
both have
\[
 E(I)=\sum_s\binom{i_s}{3}-\sum_{s<t}i_si_t=320.
\]
For the first, \(120+364-(10+14+140)=320\); for the second,
\(1+35+455-(21+45+105)=320\). Thus partitions cannot count
distinct spectral values. Under the corrected multiplicity rule their
contributions are four and eight, and must be added.
\end{example}

\begin{remark}[Exact finite checks]\oursComputed{}
\texttt{code/check\_laplacian.py} constructs the Chevalley--Eilenberg
boundary with its alternating signs, checks \(d^2=0\), and factors exact
integer characteristic polynomials for every degree of weights 1--22 and
25. A separately implemented gap-partition enumeration gives exactly the
same eigenvalues and total multiplicities, including multiplicity 12 at
eigenvalue 320 in weight 25. It also expands
\[
 \prod_{i<j}(z_i-z_j)^3\prod_i z_i^a,\qquad a=1,2,
\]
as alternating polynomials and checks both boundary and adjoint boundary
vanish for exterior degrees 1--4. No numerical eigenvalue rounding is
used. The second check, partition enumeration, shares no boundary-matrix
code with the first. These computations confirm a finite range and the
explicit correction; they do not replace the all-weight published theorem.
An independent reviewer also built the weight-25 boundaries with a different
sign algorithm and computed exact ranks of \(\Delta-320I\). The nullities
in exterior degrees 3, 4, 5 and 6 are respectively 2, 5, 4 and 1, totaling
12. Its separate enumeration checks every weight through 25, including
23 and 24, before identifying the first collision. The review code and
output are included.
\end{remark}

\subsection{Verification annotation: the literal KP normalization}
\label{dev:1977-kp-check}
Items 1977-434.3--1977-434.5 belong to the published rational-KP thread
\citep[Theorem~3 and its corollaries]{krichever1978rational}. The following
calculation makes the editorial choice of signs and momenta explicit. It
does not claim an error in a published theorem stated with other conventions.

\begin{proposition}[Pole equations in the notebook's convention]\ours{}
Let the functions \(q_i(y,t)\), \(1\leq i\leq N\), be smooth and pairwise
distinct on an open parameter domain, and put
\[
 u(x,y,t)=-2\sum_{i=1}^N(x-q_i(y,t))^{-2}.
\]
As an identity of meromorphic functions of \(x\),
\[
 3u_{yy}+(4u_t+6uu_x+u_{xxx})_x=0
\]
holds if and only if, for every \(i\),
\[
 q_{i,yy}=-8\sum_{j\ne i}(q_i-q_j)^{-3},\qquad
 q_{i,t}=\frac34q_{i,y}^{\,2}-3\sum_{j\ne i}(q_i-q_j)^{-2}.
 \tag{K1}
\]
\end{proposition}
\begin{proof}
At a fixed pole write \(s=x-q_i\) and
\(u=-2s^{-2}+R_0+R_1s+O(s^2)\), where
\[
 R_0=-2\sum_{j\ne i}(q_i-q_j)^{-2},\qquad
 R_1=4\sum_{j\ne i}(q_i-q_j)^{-3}.
\]
The leading sixth-order terms cancel. The only remaining negative powers
in the differential equation have coefficients
\[
 [s^{-4}]=48q_{i,t}-72R_0-36q_{i,y}^{\,2},\qquad
 [s^{-3}]=-24R_1-12q_{i,yy}.
\]
Their vanishing is exactly (K1). Conversely, if (K1) holds, the left-hand
side is a rational function of \(x\) with no finite poles. It tends to
zero at infinity, so it is identically zero. The distinct-pole hypothesis
is used in the expansion and is not omitted at collisions.
\end{proof}

\begin{proposition}[Compatible canonical Hamiltonians]\ours{}
With canonical coordinates \((q_i,p_i)\) on the distinct-position locus,
define
\[
 H=\frac12\sum_i p_i^2-4\sum_{i<j}(q_i-q_j)^{-2},\qquad
 K=\frac14\sum_i p_i^3-3\sum_{i<j}(p_i+p_j)(q_i-q_j)^{-2}.
 \tag{K2}
\]
Then \(\{H,K\}=0\). Taking the \(H\)-flow as \(y\) and the \(K\)-flow
as \(t\) gives \(p_i=q_{i,y}\) and (K1), hence the displayed rational KP
solution locally away from collisions.
\end{proposition}
\begin{proof}
Write \(d_{ij}=q_i-q_j\). Direct differentiation in
\(\{H,K\}=\sum_i(H_{q_i}K_{p_i}-H_{p_i}K_{q_i})\) cancels the
\(p_i^2\) terms. The remaining momentum terms are
\(-6\sum_{i\ne j}p_ip_jd_{ij}^{-3}=0\), by exchanging \(i,j\).
What remains is
\[
 -24\sum_i\sum_{j\ne i}\sum_{k\ne i}d_{ij}^{-3}d_{ik}^{-2}.
\]
Terms with \(j=k\) cancel in pairs. Group the others by unordered triples
and set \(a=q_i-q_j\), \(b=q_j-q_k\), \(c=q_k-q_i\), so
\(a+b+c=0\). Multiplication of a triple's sum by \(a^3b^3c^3\) gives
\[
 a^3(b-c)+b^3(c-a)+c^3(a-b)=0.
\]
This last polynomial vanishes upon substituting \(c=-a-b\). Thus the
bracket is zero. Hamilton's equations give
\(q_{i,y}=p_i\), \(p_{i,y}=-8\sum_{j\ne i}d_{ij}^{-3}\), and the
second formula in (K1) for \(q_{i,t}\). Commuting smooth vector fields have
local compatible flows on this open locus; no global continuation through
a collision is claimed.
\end{proof}

\begin{remark}[The lower coefficient in the third-order operator]\oursComputed{}
Let \(D=\partial_x\), \(L=D^2+u\) and
\(M=D^3+\frac32uD+w\). The coefficient of \(D\) in
\([\partial_y-L,\partial_t+M]\) vanishes precisely when
\(w_x=\frac34(u_{xx}+u_y)\). The remaining scalar equation is
\[
 u_t+w_y-w_{xx}+u_{xxx}+\frac32uu_x=0.
\]
Differentiating it in \(x\) and eliminating derivatives of \(w\) gives
the notebook's KP equation. Locally one may choose the function of
\((y,t)\) arising on integrating \(w_x\) to satisfy this scalar equation;
global boundary or periodicity conditions must be supplied separately.
Using \(\partial_t-M\) instead reverses the KP time convention. Thus the
final coefficient in \(M\) cannot in general be the same \(u\) as in \(L\).
\end{remark}

The Laurent expansion was derived independently by the record editor and
the chapter editor. \texttt{code/check\_kp.py} reproduces it symbolically,
constructs differential-operator composition by the Leibniz rule, and
checks the coefficient elimination. The triple cancellation in the
Hamiltonian proof is checked as a polynomial identity, with independent
exact evaluations of the full Poisson bracket. These checks found the
same normalization. They do not assert a classification beyond the
explicit pole ansatz and local distinct-position domain just proved.


% No new C/D theorem is claimed for records419--427.
\subsection{An annotation on the thirteen-coin convention}

This is a check of an expository puzzle (1977-424.2), not a proposed new
solution to a class C or D research question.  The classical subject appears
in \citet{dyson1946pennies}; the publisher record was checked, but the proof
below is a self-contained reconstruction.  The difference between identifying
a counterfeit and determining its sign is essential to the notebook's number.

\begin{proposition}[Identity only, with equal pan counts]\ours{}
Suppose exactly one of \(m\) coins has a different mass, either heavier or
lighter, and all other coins have a common mass.  No genuine coin is supplied
externally.  Every weighing compares equal numbers of coins on an ideal
three-outcome balance.  For \(N\geq1\), the largest \(m\) whose counterfeit
identity can always be determined in at most \(N\) weighings is
\((3^N-1)/2\).  Its sign need not always be determined.
\end{proposition}

\begin{proof}
For the upper bound let the first weighing compare \(a\) coins against
\(a\), leaving \(b=m-2a\) aside.  Either unbalanced branch contains
\(2a\) possible signed counterfeits.  They must be distinguished by
\(N-1\) subsequent ternary observations, so
\(2a\leq3^{N-1}\).  Since the left side is even,
\(2a\leq3^{N-1}-1\).  On the balanced branch the \(b\) unweighed coins
have unknown signs.  For any fixed coin, its two signs can follow the same
history only if that coin is never weighed on that history; until its first
weighing both signs produce balance, whereas its first appearance produces
opposite imbalances.  Consequently only the all-balance history can fail to
determine the sign, and it can represent at most one identity.  At least
\(2b-1\) terminal histories are therefore necessary when \(b>0\).
Hence \(b\leq(3^{N-1}+1)/2\); the same bound is trivial if \(b=0\).
Adding gives \(m\leq(3^N-1)/2\).  The argument permits adaptive choices
of the subsequent weighings.

For attainment, use the alphabet \(\{-1,0,1\}\), in cyclic order
\(-1\to0\to1\to-1\).  For every nonconstant word of length \(N\),
find its first unequal adjacent pair.  Retain the word precisely when that
pair follows the cyclic order.  Negation reverses the order, so exactly one
of each opposite pair is retained.  There are
\(M=(3^N-3)/2\) retained words.  Cyclically permuting all three alphabet
symbols preserves the selection rule, so each coordinate takes each symbol
equally often among the retained words.  Their coordinate sums are therefore
zero.  Add the all-zero word, giving \(M+1=(3^N-1)/2\) columns.

Assign one column to each coin.  In weighing \(j\), put entries \(-1\)
on the left, \(1\) on the right, and leave entries \(0\) aside.  Every
weighing has equal pan counts.  A heavy counterfeit produces its column as
the outcome word, with right-heavy denoted \(1\); a light one produces the
negative column.  Distinct nonzero columns are never opposites, so these
histories identify their coins and signs.  The all-zero history identifies
the last coin, with either sign.  For \(N=1\) the retained set is empty and
the single coin is identified without weighing, as required.
\end{proof}

\begin{remark}[If the sign is required]\ours{}
When both identity and sign must be recovered, the balanced first branch
must distinguish all \(2b\) states, hence
\(b\leq(3^{N-1}-1)/2\).  Together with the same unbalanced-branch bound
this gives \(m\leq(3^N-3)/2\).  The nonzero retained columns above attain
that bound for \(N\geq2\).  Thus three weighings permit 12 coins with
both outputs, or 13 with identity alone.  This distinction, not a numerical
misprint, accounts for the notebook's 13.
\end{remark}

\begin{example}[Exhaustive finite checks]\oursComputed{}
The accompanying \texttt{coin\_check.py} generates the columns and computes
actual left and right pan masses, independently of the signed-column
formula.  For \(N=1,\ldots,6\), it checks respectively
\(1,4,13,40,121,364\) coins.  Every coordinate balances, and every outcome
determines a unique coin.  The only two-state collision is the sign of the
zero-column coin.  In the three-weighing case all 26 signed hypotheses were
enumerated, yielding 25 histories and no collision between distinct
identities.  The proof establishes the general statement; finite enumeration
only checks the implementation and the notebook's small case.
\end{example}


% No independent C/D result is claimed by this record editor.
% Items 434.3--434.5 contain normalized KP verification annotations; the assembler independently checks them.
% Item 436.5 remains an annotated partial result: breathers are already explicitly present in the notebook.
% The chapter assembler supplies any selected deep development separately.


\subsection{Independent checks of the spectral and algebraic formulas}

The following computations explain emendations in occurrences
1977-442.4, 1977-442.6--.7, 1977-443.9 and 1977-443.12--.14, and two later
algebraic formulas. They are checks made for this companion, with no claim
of historical novelty. The relevant published subjects are described by
\citet{calogerodegasperis1976inverse,calogerodegasperis1978extension,calogero1978motion}; only their publisher metadata and abstracts or
editorial footnotes were available for the present comparison. Thus the
proofs below verify the displayed notebook equations independently of a
claim to have read their full published normalization conventions.

\begin{proposition}[Transmission phase in the notebook convention]\oursComputed{}
Let \(u\in C(\mathbb R)\cap L^1(\mathbb R)\) be real. Use the
notebook's physical scattering normalization for
\(\psi''=(u-k^2)\psi\): at negative infinity,
\(\psi\sim T(k)e^{-ikx}\), and at positive infinity,
\(\psi\sim e^{-ikx}+R(k)e^{ikx}\). With
\(T=|T|e^{2i\varphi}\) and the branch \(\varphi\to 0\) as
\(k\to+\infty\), its leading high-energy phase is
\[
 \varphi(k)=-\frac{1}{4k}\int_{\mathbb R}u(x)\,dx+o(k^{-1}).
\]
Consequently, the source's coefficient \(-c_0/2\), where
\(c_0=\int u\), is incompatible with its own factor 2 in the phase of T.
\end{proposition}
\begin{proof}
Normalize the left Jost solution f by \(f\sim e^{-ikx}\) at negative
infinity. Its Volterra equation is
\[
 f(x)=e^{-ikx}+\frac1k\int_{-\infty}^x
          \sin(k(x-y))u(y)f(y)\,dy.
\]
For real \(k>0\), Gronwall's inequality gives
\(\sup_x|f(x)|\le\exp(\|u\|_1/k)\) and
\(\sup_x|f(x)-e^{-ikx}|\le
\|u\|_1\exp(\|u\|_1/k)/k\).
Splitting the sine into exponentials shows that the coefficient of
\(e^{-ikx}\) at positive infinity is
\[
 A(k)=1-\frac1{2ik}\int_{\mathbb R}e^{iky}u(y)f(y)\,dy
     =1+\frac{i c_0}{2k}+O(k^{-2}).
\]
The error bound follows from the uniform estimate and integrability
of u. For all sufficiently large k, \(A(k)\ne0\).
The physical solution is \(\psi=f/A\), so \(T=A^{-1}\), and
\(\log T=-i c_0/(2k)+O(k^{-2})\). Taking half the imaginary part
on the branch near1 proves the formula. No convergence of a full
asymptotic series is asserted.

An independent exact check uses
\(u=-2r^2\operatorname{sech}^2(rx)\), \(r>0\), and
\[
 \psi=e^{-ikx}\frac{ik+r\tanh(rx)}{ik+r},\qquad
 T(k)=\frac{k+ir}{k-ir}.
\]
Direct differentiation gives \(\psi''=(u-k^2)\psi\). Its multiplying
factor tends to 1 at positive infinity and to
\((ik-r)/(ik+r)=(k+ir)/(k-ir)=T(k)\) at negative infinity,
so it has exactly the stated physical normalization, with \(R=0\).
Here \(c_0=-4r\) and \(\varphi=\arctan(r/k)=r/k+O(k^{-3})\),
again giving \(-c_0/4\). The symbolic residual and expansion are
recorded in the accompanying verification output. An independent
review caught an earlier reciprocal error in the stated boundary
normalization; these explicit endpoint checks correct it without
altering the final coefficient \(-c_0/4\).
\end{proof}

\begin{proposition}[Weights of the integrated recursion moments]\oursComputed{}
Let \(u\in\mathcal S(\mathbb R)\), and for
\(f\in\mathcal S(\mathbb R)\) define
\[
 L_u f=f''-4uf+2u'\int_x^\infty f(s)\,ds,\qquad
 S_u=xu'+2u,\qquad c_n[u]=\int_{\mathbb R}L_u^n S_u\,dx,
\]
The upper-tail primitive of a Schwartz function is bounded; its
product with \(u'\) is Schwartz. Thus L preserves this space,
all the displayed iterates and integrals exist, and the boundary
terms used below vanish.
For \(u_\lambda(x)=\lambda^2u(\lambda x)\), \(\lambda>0\),
\[
 c_n[u_\lambda]=\lambda^{2n+1}c_n[u].
\]
Thus a pure scaling flow \(u_t=b(xu_x+2u)\), for constant \(b\in\mathbb R\),
has \(\dot c_n=(2n+1)bc_n\), rather than \(nbc_n\).
\end{proposition}
\begin{proof}
If \(h_\lambda(x)=\lambda^a h(\lambda x)\), then its upper-tail
primitive has weight \(a-1\). Each of the three terms in
\(L_{u_\lambda}h_\lambda\) therefore has weight \(a+2\). Since
\(S_{u_\lambda}(x)=\lambda^2S_u(\lambda x)\), induction gives
weight \(2n+2\) before integration and \(2n+1\) afterwards.
The scaling-flow solution is obtained by \(\lambda=e^{bt}\).

Independently, integration by parts gives \(c_0=\int u\).
Put \(U(x)=\int_x^\infty u(s)\,ds\). Then
\(\int_x^\infty S_u=U-xu\), and the formula for L gives
\(c_1=-3\int u^2\). For example, on
\(u_\lambda=\lambda^2e^{-\lambda^2x^2}\), the two quantities are
\(\lambda\sqrt\pi\) and \(-3\lambda^3\sqrt{\pi/2}\).
These exact integrals independently test weights 1 and 3, already
refuting the source's proposed weights 0 and 1.
\end{proof}

\begin{proposition}[The sign in the Bianchi quadrilateral]\oursComputed{}
Let \(p,q\in\mathbb C\) be constants, and let U,A,B be
complex-valued \(C^1\) functions on a real interval I. Suppose the
two lower B\"acklund edges satisfy
\[
 U_x+A_x=-\tfrac12(A-U)(4p+A-U),
\]
and the same equation with \((A,p)\) replaced by \((B,q)\).
On any subinterval where the displayed denominator is nonzero,
the compatible upper value is the \(C^1\) function
\[
 C=U-\frac{(p+q)(A-B)}{p-q+\tfrac12(A-B)}.
\]
In particular, the denominator uses \(A-B\), with the same order
as the numerator in these conventions.
\end{proposition}
\begin{proof}
The upper edges require
\[
 A_x+C_x=-\tfrac12(C-A)(4q+C-A),\qquad
 B_x+C_x=-\tfrac12(C-B)(4p+C-B).
\]
Subtracting them and eliminating \(A_x-B_x\) by the two lower
edges gives a linear equation for C. Solving it yields the displayed
formula. This elimination is one check.

For a second check, differentiate that rational formula using the two
lower equations, leaving \(U_x\) arbitrary. Substitution into either
upper equation gives zero identically. The symbolic verification file
performs these as two separate residual calculations. The computation
only concerns this local differential quadrilateral; it neither fills
in missing scattering norming data nor proves a global permutability
statement at zeros of the denominator.
\end{proof}

\begin{proposition}[The Hermite interaction matrix]\oursComputed{}
Let \(\bar x_1,\ldots,\bar x_n\) be the distinct zeros of the
physicists' Hermite polynomial \(H_n\), and put
\[
 A_{jk}=-\frac1{(\bar x_j-\bar x_k)^2}\quad(j\ne k),\qquad
 A_{jj}=\sum_{k\ne j}\frac1{(\bar x_j-\bar x_k)^2}.
\]
For the source's root dynamics
\[
 \ddot x_j=-x_j+\sum_{k\ne j}
       \frac{1+2\dot x_j\dot x_k}{x_j-x_k},
\]
the equilibrium linearization is
\(\ddot\varepsilon+(I+A)\varepsilon=0\).
Moreover, \(\operatorname{Spec}(A)=\{0,1,\ldots,n-1\}\), with
each eigenvalue simple.
\end{proposition}
\begin{proof}
The derivative of each inverse difference contributes
\(-(\varepsilon_j-\varepsilon_k)/(\bar x_j-\bar x_k)^2\).
The velocity products have zero linear part at a stationary point.
The independent term \(-x_j\) contributes \(-\varepsilon_j\),
giving the stated matrix \(I+A\).

For the spectrum, use the notebook polynomial equation
\[
 \Psi_{tt}+\tfrac12\Psi_{xx}-x\Psi_x+n\Psi=0,
 \qquad
 \Psi=2^{-n}\sum_{m=0}^n c_m(t)H_{n-m}(x),\quad c_0=1.
\]
The Hermite differential equation gives
\(\ddot c_m+mc_m=0\) for \(1\le m\le n\).
At a polynomial with distinct roots, the differential of the
root-to-coefficient map is invertible: a variation has
\(\delta\Psi(\bar x_j)=-\Psi'(\bar x_j)\delta x_j\), and
evaluation at n distinct points is invertible on polynomials of
degree at most \(n-1\). Therefore \(I+A\) is conjugate to
\(\operatorname{diag}(1,\ldots,n)\).

As a second explicit check, for n=3 the roots are
\(-\sqrt{3/2},0,\sqrt{3/2}\) and
\[
 A=\begin{pmatrix}
 5/6&-2/3&-1/6\\
 -2/3&4/3&-2/3\\
 -1/6&-2/3&5/6
 \end{pmatrix},\qquad
 \det(\lambda I-A)=\lambda(\lambda-1)(\lambda-2).
\]
The verification code independently differentiates the full nonlinear
force and obtains its Jacobian \(-I-A\). It also checks the Hermite
mode identity exactly for all 36 pairs \(1\le m\le n\le8\).
Already n=1 shows why the identity cannot be omitted: A is zero but
the single root satisfies \(\ddot x=-x\).
\end{proof}

\begin{remark}[Two later algebraic checks]\oursComputed{}
For occurrence 1977-457.18, take the ordinary \(\mathfrak{sl}_2\)
relations and define
\[
 \varepsilon_{3r}=\tfrac12ht^r,\quad
 \varepsilon_{3r+1}=et^r,\quad
 \varepsilon_{3r-1}=\tfrac12ft^r.
\]
Matrix commutators verify all nine residue cases of
\([\varepsilon_i,\varepsilon_j]=a_{ij}\varepsilon_{i+j}\),
where \(a_{ij}\in\{-1,0,1\}\) is congruent to \(j-i\) modulo 3.
Separately, the 27 residue triples satisfy Jacobi directly. The
nonhomogeneous ideal \((t-1)\mathfrak{sl}_2[t,t^{-1}]\) explains why
the classification convention is graded simplicity, as already stated
in \citet[Introduction]{kac1968graded}.

For occurrence 1977-457.25, left Grassmann differentiation satisfies
\(\{\partial_{\theta^a},\theta^b\}=\delta_a^b\), while distinct
odd multiplications or derivatives anticommute. Hence the two cross
terms in the notebook's Q and \(\bar Q\) each contribute
\(-i\sigma^m\partial_m\), giving
\(\{Q_\alpha,\bar Q_{\dot\beta}\}=-2i\sigma^m_{\alpha\dot\beta}
\partial_m\); the same-chirality brackets vanish. Independently,
the code realizes multiplication and differentiation as exact 16-by-16
matrices on the exterior algebra of four generators and checks all 12
anticommutator pairs with four independent commuting derivative
symbols. These calculations validate the surviving manuscript signs;
they do not resolve the illegible cancelled mark or supply the missing
Lorentz action on the spinor coordinates.
\end{remark}
